\section{Reproducibility entry points}
\label{app:repro}
\begin{raggedright}
All statements in this paper trace to artifacts in the software repository \cite{EngiProof2026}. The entry points are:
\begin{itemize}[leftmargin=1.5em]
  \item \texttt{engiproof verify <ID>} and \texttt{engiproof verify-all} run non-mutating verification against the frozen evidence.
  \item \texttt{engiproof environment-record} verifies every study and emits a compact record of the environment and the comparator classes.
  \item \texttt{engiproof graph-audit <ID>}, \texttt{discrepancy-audit <ID>}, \texttt{discrepancy-gate <ID>} and \texttt{promotion-gate <ID>} run the provenance and gate audits.
  \item \texttt{engiproof discrepancy-taxonomy} and \texttt{ingestion-summary-audit} report the taxonomy and ingestion records.
  \item Continuous integration runs Linux, Windows and macOS with a tracked-file non-mutation gate (\texttt{.github/workflows/ci.yml}).
  \item \texttt{docs/MUTATION\_INVENTORY\_v0.2.0.md} records the pre-fix inventory, and \texttt{docs/CROSS\_ENVIRONMENT\_VERIFICATION.md} records the cross-environment runs.
\end{itemize}
Before submission, the following are to be regenerated from a tagged release and checked:
\begin{itemize}[leftmargin=1.5em]
  \item Tables~\ref{tab:crossenv} and~\ref{tab:evidence}, using \texttt{scripts/build\_tables.py};
  \item \ref{app:claims}, using \texttt{scripts/build\_claim\_appendix.py};
  \item the numbers stated in the text, using \texttt{scripts/check\_manuscript\_facts.py}.
\end{itemize}
Tables~\ref{tab:reproduction}--\ref{tab:mutation}, \ref{tab:cases} and~\ref{tab:failures} are written from the repository records and were checked in the final source audit.
Copyrighted source PDFs are not redistributed; each is identified by its SHA-256 and bibliographic record.
\end{raggedright}
